\section{Conclusion}
\label{sec:conclusion}

This paper studied residual connectivity for graph classification as a
structured branch-by-layer design problem. Across six datasets and four
message-passing operators, residual topology matters, but its value depends on
the dataset, operator, and branch budget. MatrixRes is the most frequent winner
with the best mean rank in the full benchmark, yet the per-dataset GCNConv
winners remain mixed, and ENZYMES has low absolute accuracy and should be read
cautiously. Branch-count ablations show that adding branches helps only within a
limited operating region: additional branches raise representational diversity
but also residual traffic, and they can preserve redundant branch similarity or
weaken optimization. The synthetic multi-class benchmark supports a conditional
reading---matrix-style reuse is unnecessary for the easiest two- and three-class
settings but becomes more competitive from four classes onward, where
MatrixResGated attains the best average and normalized accuracy and MatrixRes
degrades least from two to eight classes. The practical implication is that
residual topology, branch count, operator choice, and residual control should be
evaluated jointly, and that candidate settings should be rerun across folds
before a sensitivity result is treated as a performance claim.
